%! Problems with Sampling — Unit 3, Lesson 3.3 — Slides (Beamer)
\documentclass[aspectratio=169, 11pt]{beamer}
\usepackage{apstats-beamer}   % provides \navyheader{} and \sectionlabel[color]{}

\begin{document}

% TITLE SLIDE (CourseName is not defined in beamer, so the name is literal)
{
\setbeamercolor{background canvas}{bg=navy}
\begin{frame}[plain]
  \vspace{2.8cm}\hspace{0.55cm}%
  \begin{minipage}{0.9\textwidth}
    {\color{goldacc}\bfseries\small LESSON 3.3}\\[0.5cm]
    {\color{white}\bfseries\LARGE Problems with Sampling}\\[1.2cm]
    {\color{skymid}\small Statistics~$\cdot$~Unit 3}
  \end{minipage}
\end{frame}
}

% GRADUAL-RELEASE deck: targets → warm-up → I DO divider → one frame per instruction row of the
% notes (Volunteers; Coverage and Response; Response Bias, which carries the crux) → WE DO (The
% Work-Hours Survey) → spoken debrief → close & START THE HOMEWORK, which is the individual
% practice. There is no group activity and no independent practice set in the notes.
\newcommand{\redink}[1]{{\color{redacc}\bfseries #1}}

% ── LEARNING TARGETS ─────────────────────────────────────────────────────────
\begin{frame}[t, plain]
  \navyheader{Today's targets}
  \sectionlabel{Lesson 3.3}
  By the end of today, I can\ldots
  \begin{itemize}\setlength{\itemsep}{4pt}
    \item identify \textbf{voluntary response bias} and the bias in a \textbf{convenience
          sample}, and say which way each pushes the estimate.
    \item tell \textbf{undercoverage bias} from \textbf{nonresponse bias}, even in a random
          sample.
    \item explain how wording or a self-reported answer creates \textbf{response bias}.
    \item explain why a larger sample reduces variability but does not fix bias.
  \end{itemize}
  \begin{block}{How today works}
    Warm-up $\to$ \textbf{notes together} $\to$ \textbf{one we work as a class} $\to$
    debrief $\to$ \textbf{you start the homework here, alone}.
  \end{block}
\end{frame}

% ── WARM-UP ──────────────────────────────────────────────────────────────────
\begin{frame}[t, plain]
  \navyheader{Warm-up}
  \sectionlabel{5 minutes $\cdot$ on your own}
  \begin{enumerate}\setlength{\itemsep}{8pt}
    \item Name the method: a computer picks 60 of 1{,}200 students; then 15 at random from each
          of four grades.
    \item Maya asks the 12 friends at her lunch table how many hours they work. Did chance
          decide? Who is left out?
    \item Which gets more ``yes''? (A) \emph{Should our school ban phones in class?} \quad
          (B) \emph{\ldots\ so that students can focus and learn?}
  \end{enumerate}
  \vspace{6pt}
  \begin{block}{Keep items 2 and 3 on the board}
    Each one is a different way a survey goes wrong. The notes name both.
  \end{block}
\end{frame}

% ── I DO — direct instruction begins ─────────────────────────────────────────
{
\setbeamercolor{background canvas}{bg=navy}
\begin{frame}[plain]
  \vspace{1.8cm}\hspace{0.8cm}%
  \begin{minipage}{0.86\textwidth}
    {\color{goldacc}\bfseries\small GUIDED NOTES \ \textbullet\ \ I DO}\\[0.7cm]
    {\color{white}\bfseries\Large Open to your Guided Notes. Fill each blank as we name it
    --- not before.}\\[0.9cm]
    {\color{skymid}\small Four terms go in the vocabulary box today: voluntary response bias,
    undercoverage bias, nonresponse bias, response bias.}
  \end{minipage}
\end{frame}
}

% ── ROW 1: VOLUNTEERS ────────────────────────────────────────────────────────
\begin{frame}[t, plain]
  \navyheader{Who chose to be in the sample?}
  \sectionlabel{notes, Volunteers $\cdot$ problems 1--3}
  \vspace{2pt}
  \begin{columns}[t]
    \begin{column}{0.48\textwidth}
      \begin{block}{Voluntary response sample}
        People \textbf{choose} themselves; those with \textbf{strong} opinions answer most, so
        the result leans their way.
      \end{block}
      \vspace{3pt}
      \begin{block}{Convenience sample}
        Whoever is \textbf{easiest} to reach. Neither one uses \textbf{chance}.
      \end{block}
    \end{column}
    \begin{column}{0.48\textwidth}
      {\small\textbf{A complete answer, every time}}
      \begin{enumerate}\setlength{\itemsep}{3pt}\small
        \item Name the bias.
        \item Who is over- or under-represented, and why --- in context.
        \item The \redink{direction}: too high or too low?
      \end{enumerate}
      \vspace{6pt}
      {\small A traffic poll? An AP Calculus class? Comment cards --- \redink{which way} do
      those lean?}
    \end{column}
  \end{columns}
\end{frame}

% ── ROW 2: COVERAGE AND RESPONSE ─────────────────────────────────────────────
\begin{frame}[t, plain]
  \navyheader{Who is missing?}
  \sectionlabel{notes, Coverage and Response $\cdot$ problems 4--6}
  \vspace{2pt}
  \begin{columns}[t]
    \begin{column}{0.48\textwidth}
      \begin{block}{Undercoverage}
        Part of the population has a \textbf{reduced} chance --- or none --- of being chosen,
        often because the \textbf{list} leaves them out.
      \end{block}
    \end{column}
    \begin{column}{0.48\textwidth}
      \begin{block}{Nonresponse}
        Individuals \textbf{chosen} for the sample can't be reached or \textbf{refuse}, and
        they may differ from those who answer.
      \end{block}
    \end{column}
  \end{columns}
  \vspace{8pt}
  \begin{block}{The test}
    Could they ever have been picked? \emph{No} $\to$ undercoverage. \quad \emph{Picked, but no
    answer} $\to$ nonresponse. \quad Either one can spoil a \redink{random} sample.
  \end{block}
\end{frame}

% ── ROW 3: RESPONSE BIAS — AND THE CRUX ──────────────────────────────────────
\begin{frame}[t, plain]
  \navyheader{How was it asked --- and does size help?}
  \sectionlabel{notes, Response Bias $\cdot$ problems 7--10 $\cdot$ the whole point of today}
  \vspace{2pt}
  \textbf{Response bias} comes from \emph{how} the data are gathered: a \textbf{leading} question,
  or an answer people won't give \textbf{honestly}.
  \vspace{6pt}
  \begin{columns}[t]
    \begin{column}{0.46\textwidth}
      \begin{block}{Game forum survey}
        \textbf{48{,}000} players chose to reply.\\
        \textbf{81\%} dislike the update.
      \end{block}
    \end{column}
    \begin{column}{0.46\textwidth}
      \begin{block}{Random sample}
        \textbf{600} players picked at random.\\
        \textbf{47\%} dislike the update.
      \end{block}
    \end{column}
  \end{columns}
  \vspace{6pt}
  ``48{,}000 is so many that 81\% must be close to the truth.'' \emph{Agree?}
  \vspace{6pt}
  \begin{block}{Say it, write it in the box}
    A larger sample reduces \redink{variability}, not \redink{bias}. More volunteers are just
    more of the same volunteers.
  \end{block}
\end{frame}

% ── WE DO ────────────────────────────────────────────────────────────────────
\begin{frame}[t, plain]
  \navyheader{Guided Practice: The Work-Hours Survey}
  \sectionlabel[goldacc]{We do $\cdot$ you hold the pen $\cdot$ problems 11--14}
  A university wants the mean hours per week its students work at paid jobs.
  \vspace{4pt}
  \begin{center}
  {\small\renewcommand{\arraystretch}{1.3}
  \begin{tabular}{@{} c l @{}}
    \textbf{Step} & \textbf{The plan} \\ \hline
    1 & 500 students picked at random from the list of students in campus dorms \\
    2 & Each is asked to meet their advisor; only 140 show up \\
    3 & The advisor asks, ``How many hours a week do you work at a paid job?'' \\
  \end{tabular}}
  \end{center}
  \begin{itemize}\setlength{\itemsep}{3pt}
    \item For each step: name the bias, say who it misses or misleads, and give the direction.
    \item ``Email all 30{,}000 and get 6{,}000 replies --- that fixes it.'' Respond.
  \end{itemize}
\end{frame}

% ── DEBRIEF ──────────────────────────────────────────────────────────────────
\begin{frame}[t, plain]
  \navyheader{Debrief: name it, aim it, and don't count on size}
  \sectionlabel{8 minutes $\cdot$ pens down, papers up --- we say these aloud}
  \vspace{6pt}
  \begin{columns}[t]
    \begin{column}{0.46\textwidth}
      \begin{block}{Big but biased}
        48{,}000 forum replies: \textbf{81\%}.\\
        6{,}000 emailed volunteers: still biased.
      \end{block}
    \end{column}
    \begin{column}{0.46\textwidth}
      \begin{block}{Small but random}
        600 at random: \textbf{47\%}.\\
        Chance gave no group a head start.
      \end{block}
    \end{column}
  \end{columns}
  \vspace{10pt}
  \begin{block}{Two things to say out loud}
    Why is ``it's biased'' not a complete answer? And why did the comment cards lean
    \redink{low} when the traffic poll leaned \redink{high}?
  \end{block}
\end{frame}

% ── CLOSE & START THE HOMEWORK ──────────────────────────────────────────────
\begin{frame}[t, plain]
  \navyheader{Homework --- start it now}
  \sectionlabel[goldacc]{5 items $\cdot$ scored 1 to 5 $\cdot$ due the first class after two study halls}
  \begin{itemize}\setlength{\itemsep}{5pt}
    \item Part A: teen drivers and phones --- a social-media poll, a trooper's random phone
          survey, and the teens who never pick up. Name each bias and its direction.
    \item Part B: one multiple-choice item, with a sentence of justification, and a reporter
          who trusts the bigger poll.
  \end{itemize}
  \begin{block}{Silent and alone --- I am circulating}
    \emph{Part B, item 2} is the one I am reading over your shoulder. The AP Practice page is
    \textbf{extra credit, scored 1 to 5} --- work it for the reps.
  \end{block}
\end{frame}

% ── CLOSE ────────────────────────────────────────────────────────────────────
{
\setbeamercolor{background canvas}{bg=navy}
\begin{frame}[plain]
  \vspace{1.5cm}\hspace{0.8cm}%
  \begin{minipage}{0.86\textwidth}
    {\color{goldacc}\bfseries\small BEFORE YOU GO}\\[0.7cm]
    {\color{white}\bfseries\Large Name the bias, say who it favors, and say which way it
    pushes. A larger sample reduces variability --- not bias.}\\[0.9cm]
    {\color{skymid}\small Next: Lesson 3.4 --- Designing Experiments: treatments, random
    assignment, and why only an experiment can show cause.}
  \end{minipage}
\end{frame}
}

\end{document}
